\section{Networks with row-stochastic coupling and main results}
\label{sec:rw-model-results}
\label{sec:model-results}

We study how a perturbation with zero stationary projection changes a local
fold-matching condition.  We first state the RFDE result, including the
parameter selected by that finite-interval condition.  The subsequent range
and norm estimates describe which transverse perturbations can produce the
response and when the bounds remain uniform in network size.

\subsection{Network, phase space, and projection}

For each \(N\geq2\), let \(P_N\) be a nonnegative row-stochastic matrix and
let \(\pi_N\) be a strictly positive stationary probability vector.  We
assume that there are constants \(D>0\) and \(\gamma>0\), independent of
\(N\), such that
\begin{equation}
 P_N\mathbf1=\mathbf1,
 \qquad \pi_N^TP_N=\pi_N^T,
 \qquad \pi_N^T\mathbf1=1,
 \qquad \tau(P_N)\leq1-\gamma .
 \label{eq:rw-markov-assumptions}
\end{equation}
Here \(\tau(P_N)\) is the Dobrushin contraction coefficient.  Define the
collective and transverse projections and the transverse generator by
\begin{equation}
 P_{c,N}=\mathbf1\pi_N^T,
 \qquad P_{\perp,N}=\Id-P_{c,N},
 \qquad E_N=\Ker\pi_N^T,
 \qquad A_N=D(P_N-\Id)|_{E_N}.
 \label{eq:rw-collective-transverse}
\end{equation}
We use the dimension-independent norm
\begin{equation}
 \norm{x}_N=\abs{\pi_N^Tx}+\osc(x),
 \qquad
 \osc(x)=\max_i x_i-\min_i x_i.
 \label{eq:rw-network-norm}
\end{equation}
In this norm the Dobrushin assumption gives
\begin{equation}
 \norm{A_N^{-1}}_{E_N\to E_N}\leq\frac1{D\gamma},
 \qquad
 \norm{A_Nz}_N\leq D(2-\gamma)\norm{z}_N
 \quad (z\in E_N).
 \label{eq:rw-dobrushin-resolvent}
\end{equation}

Let \(c_N=(c_{1,N},\ldots,c_{N,N})^T\) be the vector of quadratic
coefficients.  We impose the uniform bounds
\begin{equation}
 0<c_-\leq c_{i,N}\leq c_+,
 \qquad \pi_N^Tc_N=\alpha>0,
 \label{eq:rw-curvature-assumptions}
\end{equation}
where \(c_-\), \(c_+\), and \(\alpha\) do not depend on \(N\).  Fix
\(\beta>0\), \(K\neq0\), and delay locations
\begin{equation}
 0\leq\theta_0<\theta_1<\cdots<\theta_m\leq\Theta_*.
 \label{eq:rw-delay-support}
\end{equation}
The matrices \(B_{k,N}\) are fixed delayed-coupling matrices satisfying
\begin{equation}
 \sup_{N\geq2}\sum_{k=0}^m
 \norm{B_{k,N}}_{\cL(\mathbb R^N,\norm{\cdot}_N)}<\infty.
 \label{eq:rw-base-layer-bound}
\end{equation}

An admissible perturbation of the delayed-coupling matrices is a tuple
\(\eta=(\Delta B_0,\ldots,\Delta B_m)\) in
\begin{align}
 \mathcal T_N
 &=\left\{\eta=(\Delta B_0,\ldots,\Delta B_m):
       \pi_N^T\Delta B_k=0\ \text{for every }k,
       \quad \sum_{k=0}^m\Delta B_k=0\right\},
 \label{eq:rw-structural-parameter-space}\\
 \norm{\eta}_{\mathcal T_N}
 &=\sum_{k=0}^m
   \norm{\Delta B_k}_{\cL(\mathbb R^N,\norm{\cdot}_N)},
 \qquad B_{k,N}(\eta)=B_{k,N}+\Delta B_k.
 \label{eq:rw-structural-parameter-norm}
\end{align}
Thus the sum is unchanged:
\(\sum_kB_{k,N}(\eta)=\sum_kB_{k,N}\).  The stronger identities
\(\pi_N^T\Delta B_k=0\) imply that the perturbation vanishes after application of
the stationary projection.
The space \(\mathcal T_N\) is an unrestricted matrix tangent space: it does
not impose a fixed zero pattern, symmetry, layerwise row sums, or sign
conditions.  It provides a baseline full-range result.  Prescribed matrix
patterns are treated separately in
\cref{thm:rw-structured-delay-range}; nonnegativity is a local inequality
constraint and is preserved for sufficiently small perturbations whenever
the relevant entries of the base layers are strictly positive.

For \(\eps=\delta^2\), \(a=1+\delta^2\nu\), and
\(\mu=a-1=\delta^2\nu\), consider
\begin{align}
 \dot v(t)={}&\frac23\mathbf1-w(t)\mathbf1
 -c_N\circ(v(t)-\mathbf1)^{\circ2}
 -\frac\beta3(v(t)-\mathbf1)^{\circ3}
 +D(P_N-\Id)v(t)\notag\\
 &+\delta^2K\sum_{k=0}^mB_{k,N}(\eta)
 \left[v(t)-v\left(t-\frac{\theta_k}{\delta}\right)\right],
 \label{eq:rw-voltage-rfde}\\
 \dot w(t)={}&\delta^2\{\pi_N^Tv(t)-a\}.
 \label{eq:rw-original-recovery}
\end{align}
The phase space at fixed \(\delta\) is
\begin{equation}
 \mathcal H_{N,\delta}
 =C([ -\theta_m/\delta,0];\mathbb R^N)\times\mathbb R.
 \label{eq:rw-phase-space}
\end{equation}
For \(\Phi=(\phi,\omega)\in\mathcal H_{N,\delta}\), we use the norm
\begin{equation}
 \norm{\Phi}_{\mathcal H_{N,\delta}}
 =\sup_{-\theta_m/\delta\leq\theta\leq0}\norm{\phi(\theta)}_N
   +\abs{\omega}.
 \label{eq:rw-phase-space-norm}
\end{equation}
Let \(\mathcal F_{N,\delta,\nu,\eta}\) denote the right-hand side of
\cref{eq:rw-voltage-rfde,eq:rw-original-recovery}, evaluated on a full
history, and let
\begin{equation}
 \Pi_N(x,y)=(\pi_N^Tx,y).
 \label{eq:rw-stationary-projection}
\end{equation}
Then, for every \(\Phi\in\mathcal H_{N,\delta}\),
\begin{equation}
 \Pi_N\mathcal F_{N,\delta,\nu,\eta}(\Phi)
 =\Pi_N\mathcal F_{N,\delta,\nu,0}(\Phi).
 \label{eq:rw-projection-identity}
\end{equation}
Indeed, left multiplication by \(\pi_N^T\) annihilates each
\(\Delta B_k\)-term separately.  Equation \eqref{eq:rw-projection-identity}
compares the two projected vector fields at the same full vector history.
It does not define a closed RFDE for \(\pi_N^Tv\): through the heterogeneous
quadratic terms, a network-mean-zero component can contribute to the
collective equation.

The coefficients entering the finite-interval fold expansion are
\begin{equation}
 z_{2,N}=A_N^{-1}P_{\perp,N}c_N,
 \qquad
 \kappa_N=\frac\beta3+2\pi_N^T\diag(c_N)z_{2,N},
 \qquad
 \nu_{0,N}=-\frac{3\kappa_N}{8\alpha^3}.
 \label{eq:rw-fold-center}
\end{equation}

\subsection{The local matching problem and its sensitivity}

In fold time \(s=\delta t\), use the coordinates
\begin{equation}
 w=\frac23-\delta^2Y,\qquad
 v=\mathbf1+\delta\mathbf1X+
       \delta^2\{z_{2,N}X^2+h\},\qquad h\in E_N.
 \label{eq:rw-fold-coordinates}
\end{equation}
The singular collective field and its distinguished complete orbit are
\begin{equation}
 q_0(X,Y)=(Y-\alpha X^2,-X),\qquad
 \gamma_0(s)=\left(-\frac{s}{2\alpha},\frac{s^2-2}{4\alpha}\right).
 \label{eq:rw-singular-canard}
\end{equation}
The local construction uses a window
\(S_\delta=\sqrt{2p_0\log(1/\delta)}\), where \(p_0>4\), with fixed
buffers for the delay translations.  On a locally invariant manifold of
compatible histories, two planar traces follow the fold passage from its
opposite ends to the phase section \(X=0\).  Their endpoint conditions are
specified in \cref{eq:finite-trace-BVP}; their signed difference in the
singular first integral defines \(D_N^{\rm fin}\).

For \(\eta=(\Delta B_0,\ldots,\Delta B_m)\in\mathcal T_N\), define
\begin{equation}
 \mathsf S_N\eta
 =\frac{K}{2\alpha}P_{\perp,N}
    \left(\sum_{k=0}^m\theta_k\Delta B_k\right)\mathbf1.
 \label{eq:rw-transverse-source}
\end{equation}
This is the first moment of the perturbation matrices with respect to the
delay values, evaluated in the network-mean-zero space.
All operator and dual norms below are induced by
\eqref{eq:rw-network-norm} and \eqref{eq:rw-structural-parameter-norm}.

The Dobrushin condition makes \(A_N\) invertible on \(E_N\).
Define the quadratic-coefficient-dependent functional
and the resulting sensitivity functional by
\begin{equation}
 \mathfrak r_N(h)=-\frac1\alpha\pi_N^T\diag(c_N)h,
 \qquad
 \Lambda_N=\mathfrak r_NA_N^{-1}\mathsf S_N.
 \label{eq:rw-return-functional-definition}
\end{equation}
Equivalently,
\begin{equation}
 \Lambda_N(\eta)
 =-\frac{K}{2\alpha^2}\pi_N^T\diag(c_N)A_N^{-1}P_{\perp,N}
   \left(\sum_{k=0}^m\theta_k\Delta B_k\right)\mathbf1.
 \label{eq:rw-return-functional}
\end{equation}
Thus \(\Lambda_N\) is a linear functional on the perturbation space; for a
fixed direction \(\eta\), the scalar \(\Lambda_N(\eta)\) is the corresponding
first-order coefficient.

\begin{theorem}[Local RFDE fold matching and its sensitivity]
\label{thm:rw-local-fold-matching}
Assume \eqref{eq:rw-markov-assumptions}--\eqref{eq:rw-base-layer-bound}.
There are \(\delta_0>0\), \(\eta_*>0\), and a compact interval \(I_\nu\)
containing all \(\nu_{0,N}\) with a common positive distance to its boundary,
independent of \(N\), such that the RFDE has a two-dimensional locally
invariant manifold of compatible histories in the fold neighborhood and the
finite boundary-value problems in \cref{eq:finite-trace-BVP} define the scalar
matching function \(D_N^{\rm fin}\) in
\cref{eq:finite-matching-definition}.  Uniformly for
\(0<\delta<\delta_0\), \(\nu\in I_\nu\), and
\(\norm{\eta}_{\mathcal T_N}\le\eta_*\),
\begin{align}
 \frac{D_N^{\rm fin}}{\delta}
 &=\sqrt{2\pi}(\nu-\nu_{0,N})+O(\delta),
                                                        \label{eq:finite-comparison-level}\\
 \partial_\nu D_N^{\rm fin}
 &=\delta\sqrt{2\pi}+O(\delta^2),
                                                        \label{eq:finite-comparison-slope}\\
 \bigl\|D_\eta D_N^{\rm fin}
       +\delta^2\sqrt{2\pi}\Lambda_N\bigr\|
 &\le C\{\delta^3+\delta^2\|\eta\|\}.             \label{eq:finite-comparison-structural}
\end{align}
The structural derivative \(D_\eta D_N^{\rm fin}\) takes values in
\(\mathcal T_N^*\); derivatives are defined on an open parameter
neighborhood containing the displayed closed set.  There is also a common
\(c_0>0\) such that, after reducing \(\delta_0\) and \(\eta_*\), the equation
\(D_N^{\rm fin}(\delta,\nu,\eta)=0\) has a unique solution
\(\nu=\nu_N^{\rm fin}(\delta,\eta)\) in
\(\abs{\nu-\nu_{0,N}}<c_0\).  Put
\(\mu_N^{\rm fin}=\delta^2\nu_N^{\rm fin}\).  For fixed \(\delta\), this
map is continuously Fr\'echet differentiable on the open \(\eta\)-ball, and
\begin{align}
 \norm{D_\eta\mu_N^{\rm fin}(\delta,\eta)-\delta^3\Lambda_N}
 &\le C\{\delta^4+\delta^3\norm{\eta}_{\mathcal T_N}\},
 \label{eq:rw-local-root-sensitivity}\\
 \abs{\mu_N^{\rm fin}(\delta,\eta)-\mu_N^{\rm fin}(\delta,0)
              -\delta^3\Lambda_N(\eta)}
 &\le C\{\delta^4\norm{\eta}_{\mathcal T_N}
            +\delta^3\norm{\eta}_{\mathcal T_N}^2\}.
 \label{eq:rw-local-root-increment}
\end{align}
The zero is selected by the specified finite boundary-value problems.
It is not asserted to be a heteroclinic connection of the original RFDE.
\end{theorem}


The projection identity compares vector fields at the same history.  The
matching condition compares histories that themselves depend on the
perturbation.  The theorem quantifies the resulting indirect contribution:
it enters \(D_N^{\rm fin}\) at order \(\delta^2\) and the physical parameter
\(\mu_N^{\rm fin}\) at order \(\delta^3\).  The construction of the history
manifold and the uniform estimates are proved in \cref{sec:fold-passage}
and Appendix~A.  The conversion from the matching estimates to its unique
local zero uses the nonvanishing \(\nu\)-derivative.

The next results characterize the linear functional in this theorem.  They
also distinguish unrestricted matrix perturbations from those with a fixed
support or prescribed row sums.

\subsection{Range and dimension-uniform bounds}

\begin{theorem}[Range and a norm-optimal right inverse for the first delay moment]
\label{thm:rw-first-moment-map}
Let \(N\ge2\), let \(\pi_N\) be a strictly positive probability vector, and
put \(E_N=\ker\pi_N^T\).  Fix \(\alpha>0\), \(K\ne0\), and at least two
distinct delay values
\[
 0\le \theta_0<\theta_1<\cdots<\theta_m,
 \qquad m\ge1.
\]
Define the norm \(\norm{\cdot}_N\), the constrained perturbation space
\(\mathcal T_N\), and \(\mathsf S_N\) by
\cref{eq:rw-network-norm,eq:rw-structural-parameter-space,%
eq:rw-structural-parameter-norm,eq:rw-transverse-source},
and write \(\Delta_\theta=\theta_m-\theta_0\).  Then
\(\mathsf S_N:\mathcal T_N\to E_N\) is surjective.  A right
inverse is
\begin{equation}
 \begin{aligned}
  (\mathsf Q_Ny)_0
   &=-\frac{2\alpha}{K\Delta_\theta}y\pi_N^T,
  & (\mathsf Q_Ny)_m
   &= \frac{2\alpha}{K\Delta_\theta}y\pi_N^T,\\
  (\mathsf Q_Ny)_k&=0
  &&(k\notin\{0,m\}),
 \end{aligned}
 \qquad y\in E_N,
 \label{eq:rw-source-right-inverse}
\end{equation}
and, with respect to the norms in
\cref{eq:rw-network-norm,eq:rw-structural-parameter-norm},
the operator norms satisfy the optimal equalities
\begin{equation}
 \norm{\mathsf S_N}
 =\frac{\abs{K}\Delta_\theta}{4\alpha},
 \qquad
 \norm{\mathsf Q_N}
 =\frac{4\alpha}{\abs{K}\Delta_\theta}.
 \label{eq:rw-source-sharp-norms}
\end{equation}
If delayed-coupling terms with the same delay are first combined and only one
delay value remains, then \(\mathsf S_N=0\) on
\(\mathcal T_N\).
\end{theorem}

The preceding right inverse uses dense rank-one matrices in general.  The
next result records exactly what remains under a prescribed linear matrix
pattern.

\begin{theorem}[Range under prescribed matrix patterns]
\label{thm:rw-structured-delay-range}
Let \(k_-\) and \(k_+\) index two delays with
\(\theta_{k_-}<\theta_{k_+}\), and let
\[
 \mathcal U_N\subset
 \{R\in\mathbb R^{N\times N}:\pi_N^TR=0\}
\]
be a linear space of allowed matrix directions.  Restrict
\(\mathsf S_N\) to tuples with
\[
 R_{k_-}=R,\qquad R_{k_+}=-R,\qquad
 R\in\mathcal U_N,
\]
and all other components zero.  Then
\begin{equation}
 \operatorname{Ran}\mathsf S_N
 =\left\{-\frac{K(\theta_{k_+}-\theta_{k_-})}{2\alpha}
           R\mathbf1:R\in\mathcal U_N\right\}.        \label{eq:rw-structured-range}
\end{equation}
In particular, the restricted map is onto \(E_N\) if and only if
\(R\mapsto R\mathbf1\) maps \(\mathcal U_N\) onto \(E_N\).  It vanishes if
every allowed direction preserves row sums.  If \(P_N\) is irreducible, the
fixed-support space containing
\[
 \{(P_N-\Id)\diag(d):d\in E_N\}
\]
satisfies the range condition.  These directions use only the support of
\(P_N-\Id\).  If both base delay layers are nonnegative and are strictly
positive on that support, sufficiently small perturbations remain
nonnegative and preserve their existing support.  The size of a structured
right inverse can depend on this support and need not be uniform in \(N\).

Under the Dobrushin hypothesis, however, the generator-supported pattern has
the right inverse
\begin{equation}
 \mathsf J_Ny=(P_N-\Id)\diag\!\left(
       \{(P_N-\Id)|_{E_N}\}^{-1}y\right),
 \qquad \mathsf J_Ny\,\mathbf1=y,                    \label{eq:rw-generator-right-inverse}
\end{equation}
with
\begin{equation}
 \norm{\mathsf J_N}\le\frac{2(2-\gamma)}{\gamma}.
 \label{eq:rw-generator-right-inverse-bound}
\end{equation}
Consequently the structured two-layer first-moment map has a right inverse
of norm at most
\begin{equation}
 \frac{8\alpha(2-\gamma)}
      {\abs{K}(\theta_{k_+}-\theta_{k_-})\gamma}.      \label{eq:rw-structured-source-inverse-bound}
\end{equation}
For example, take
\(B_{k_-}=B_{k_+}=(P_N+\Id)/2\).  Along
\(R=(P_N-\Id)\diag(d)\), the layers
\(B_{k_-}+\zeta R\) and \(B_{k_+}-\zeta R\) remain nonnegative with the same
support whenever \(\abs\zeta\norm{d}_\infty<1/2\).  They do not preserve the
row sum of each layer separately; imposing that stronger condition would
make \(\mathsf S_N\) vanish.
\end{theorem}

\begin{corollary}[Dimension-uniform Fredholm sensitivity functional for the
fold problem]
\label[corollary]{cor:rw-fold-fredholm-coefficient}
Assume \eqref{eq:rw-markov-assumptions}--\eqref{eq:rw-base-layer-bound} and
the distinct-delay hypothesis of \cref{thm:rw-first-moment-map}.
Linearization at \(\eta=0\)
of the local fold problem in \cref{sec:fold-passage}, followed by elimination
of the network-transverse equation and pairing with the one-dimensional
collective cokernel, gives the bounded linear functional
\(\Lambda_N=\mathfrak r_NA_N^{-1}\mathsf S_N\) from
\eqref{eq:rw-return-functional}.  Its norm satisfies
\begin{equation}
 \abs{\Lambda_N(\eta)}
 \leq\frac{\abs{K}c_+\Theta_*}{2\alpha^2D\gamma}
        \norm{\eta}_{\mathcal T_N}.
 \label{eq:rw-functional-uniform-bound}
\end{equation}
More generally, for any \(\mathfrak r\in E_N^*\), let
\(\Lambda_{\mathfrak r}=\mathfrak r A_N^{-1}\mathsf S_N\).  Then
\begin{equation}
 \mathfrak r(z)=\Lambda_{\mathfrak r}(\mathsf Q_NA_Nz),
 \qquad z\in E_N,
 \label{eq:rw-dual-recovery}
\end{equation}
and
\begin{equation}
 \frac{\abs{K}\Delta_\theta}{4\alpha D(2-\gamma)}
       \norm{\mathfrak r}
 \leq\norm{\Lambda_{\mathfrak r}}
 \leq
 \frac{\abs{K}\Delta_\theta}{4\alpha D\gamma}
       \norm{\mathfrak r}.
 \label{eq:rw-dual-condition}
\end{equation}
Thus the condition number of
\(\mathfrak r\mapsto\Lambda_{\mathfrak r}\), as a map onto its range,
is at most \((2-\gamma)/\gamma\), uniformly in \(N\).
\end{corollary}

The dimension-uniform nonzero response is realized by concrete growing
networks and is not a
formal consequence of writing constants independently of \(N\).

\begin{proposition}[Two growing network families]
\label[proposition]{prop:rw-growing-families}
For \(N\ge2\), let
\begin{equation}
 \pi_N=N^{-1}\mathbf1,\qquad
 P_N^{\rm mix}=(1-\rho)\Id+\rho\mathbf1\pi_N^T,\qquad
 0<\rho\le1,                                          \label{eq:rw-mixing-family}
\end{equation}
and define
\begin{equation}
 q_{i,N}=\sqrt{\frac{12}{N^2-1}}
       \left(i-\frac{N+1}{2}\right),\qquad
 c_N=\mathbf1+\sigma q_N,\qquad 0<\sigma<\frac1{\sqrt3}.
 \label{eq:rw-growing-curvature}
\end{equation}
For delays \((\theta_0,\theta_1)=(1,2)\) and the direction
\(R_0=q_N\pi_N^T\), \(R_1=-q_N\pi_N^T\), one has
\begin{equation}
 \tau(P_N^{\rm mix})=1-\rho,\qquad
 \Lambda_N(R_0,R_1)=-\frac{K\sigma}{2D\rho}           \label{eq:rw-growing-response}
\end{equation}
for every \(N\).  The coefficients \(c_{i,N}\) are pairwise distinct and
remain in a common positive interval, while
\begin{equation}
 \norm{(R_0,R_1)}_{\mathcal T_N}
 =2\norm{q_N\pi_N^T}<4\sqrt3                         \label{eq:rw-growing-direction-bound}
\end{equation}
uniformly in \(N\).  Thus the nonzero value is not obtained by allowing the
perturbation direction to grow with network size.

In contrast, let \(N=2n\ge4\) and
\(P_N^{\rm cyc}=(1-\rho)\Id+\rho C_N\), where \(C_N\) is the cyclic shift.
Then \(\tau(P_N^{\rm cyc})=1\), and on \(E_N\)
\begin{equation}
 \left\|\{D(P_N^{\rm cyc}-\Id)|_{E_N}\}^{-1}\right\|
 \ge \frac{N}{4D\rho}.                               \label{eq:rw-cycle-degeneration}
\end{equation}
Thus the dimension-uniform conclusion covers uniformly mixing families but
not general sequences of sparse local networks.
\end{proposition}


The range theorems and growing-family calculations are proved in
\cref{sec:fredholm-sensitivity}.  The same section gives the collective
cokernel, the block solvability formula, and the two-sided norm bounds for
\(\mathfrak r\mapsto\mathfrak r A_N^{-1}\mathsf S_N\).
